\documentclass[11pt]{article}
\usepackage[letterpaper,landscape,margin=0.65in]{geometry}
\usepackage{amsmath}
\usepackage{pgffor,etoolbox}

% Answer-box macros: {rows}{columns}{point size}{phantom text}.
% Augmented boxes require at least two columns.
% Construct an ordinary array; its caller decides what to hide.
\newcommand{\AnswerArray}[2]{%
	\gdef\AnswerBody{}%
	\edef\AnswerColumnCount{\number\numexpr#2\relax}%
	\foreach \row in {1,...,#1} {%
		\foreach \col in {1,...,\AnswerColumnCount} {%
			\gappto\AnswerBody{\AnswerCell}%
			\ifnum\col<\AnswerColumnCount
			\gappto\AnswerBody{&}%
			\else
			\gappto\AnswerBody{\\}%
			\fi
		}%
	}%
	\begin{array}{*{#2}{c}}\AnswerBody\end{array}%
}
\newcommand{\AnswerBox}[5][plain]{%
	\begingroup
	\def\AnswerCell{\hbox{\fontsize{#4}{#4}\selectfont #5}}%
	\def\AnswerStyle{#1}\def\PlainStyle{plain}\def\AugStyle{augmented}%
	\setlength{\arraycolsep}{2pt}%
	\renewcommand{\arraystretch}{1.4}%
	\ensuremath{\boxed{%
			\ifx\AnswerStyle\PlainStyle\else\left(\fi
			\ifx\AnswerStyle\AugStyle
			% Hide the two blocks separately so the dividing rule stays visible.
			\begin{array}{@{}c|c@{}}
				\phantom{\AnswerArray{#2}{\numexpr#3-1\relax}} &
				\phantom{\AnswerArray{#2}{1}}
			\end{array}%
			\else
			\phantom{\AnswerArray{#2}{#3}}%
			\fi
			\ifx\AnswerStyle\PlainStyle\else\right)\fi
	}}%
	\endgroup
}
\newcommand{\MatrixAnswerBox}[4]{\AnswerBox[matrix]{#1}{#2}{#3}{#4}}
\newcommand{\AugmentedAnswerBox}[4]{\AnswerBox[augmented]{#1}{#2}{#3}{#4}}

\pagestyle{empty}
\begin{document}
	\section*{Quiz answer boxes: phantom arrays}
	Every entry is an invisible \texttt{MMM} set at 24\,pt.
	Arguments are \texttt{\{rows\}\{columns\}\{pts\}\{letters\}}; the augmented column count includes the right-hand column.
	
	\bigskip
	\begin{tabular}{@{}lccc@{}}
		& Plain & Parentheses & Augmented\\[6pt]
		$1\times2$ & \AnswerBox{1}{2}{24}{MMM} & \MatrixAnswerBox{1}{2}{24}{MMM} & \AugmentedAnswerBox{1}{2}{24}{MMM}\\[14pt]
		$2\times3$ & \AnswerBox{2}{3}{24}{MMM} & \MatrixAnswerBox{2}{3}{24}{MMM} & \AugmentedAnswerBox{2}{3}{24}{MMM}\\[14pt]
		$3\times3$ & \AnswerBox{3}{3}{24}{MMM} & \MatrixAnswerBox{3}{3}{24}{MMM} & \AugmentedAnswerBox{3}{3}{24}{MMM}
	\end{tabular}
	
	\bigskip\noindent
	\verb|\AnswerBox{2}{3}{24}{MMM}|\par
	\verb|\MatrixAnswerBox{2}{3}{24}{MMM}|\par
	\verb|\AugmentedAnswerBox{2}{3}{24}{MMM}|
	
	\medskip\noindent
	The outer box uses the standard \verb|\boxed| command from \texttt{amsmath},
	with its default gap (\verb|\fboxsep|, normally 3\,pt). Parentheses and the augmentation rule are standard
	LaTeX array notation. Sizes now follow the 24\,pt \texttt{MMM} block rather
	than a prescribed golden ratio. All macro definitions are included at the top of this file.
\end{document}